%%%%%
%%%%%  Naudokite LUALATEX, ne LATEX.
%%%%%
%%%%
\documentclass[]{VUMIFTemplateClass}

\usepackage{indentfirst}
\usepackage{amsmath, amsthm, amssymb, amsfonts}
\usepackage{mathtools}
\usepackage{physics}
\usepackage{graphicx}
\usepackage{verbatim}
\usepackage[hidelinks]{hyperref}
\usepackage{color,algorithm,algorithmic}
\usepackage[nottoc]{tocbibind}
\usepackage{tocloft}
\usepackage[version=4]{mhchem}
\usepackage{textgreek}
\usepackage{textcomp}
\usepackage{float}
\usepackage{booktabs}
\usepackage{listings}
\lstset{
    language=Python,
    basicstyle=\ttfamily\footnotesize,
    keywordstyle=\bfseries,
    commentstyle=\itshape,
    stringstyle=\ttfamily,
    breaklines=true,
    frame=single,
    showstringspaces=false,
    numbers=left,
    numberstyle=\tiny,
    xleftmargin=1.5em,
}

\usepackage{titlesec}
\newcommand{\sectionbreak}{\clearpage}

\makeatletter
\renewcommand{\fnum@algorithm}{\thealgorithm}
\makeatother
\renewcommand\thealgorithm{\arabic{algorithm} algoritmas}

\usepackage{biblatex}
\bibliography{bibliografija}

\newcommand{\EE}{\mathbb{E}\,}
\newcommand{\ee}{{\mathrm e}}
\newcommand{\RR}{\mathbb{R}}

\studijuprograma{Informatikos}
\darbotipas{Ataskaita}
\darbopavadinimas{Kietųjų dalelių koncentracijos prognozavimas laiko eilutėse}
\darbopavadinimasantras{Bazinio KNN metodo su svoriais rezultatai}
\autorius{Mindaugas Žiemys}

\vadovas{Linas Litvinas, Asist., Dr.}

\begin{document}
    \selectlanguage{lithuanian}

    \onehalfspacing
    \input{titulinis}

    \singlespacing

    \tableofcontents
    \onehalfspacing


    \section{Eksperimento konfigūracija}

    Bazinė KNN metodo su svoriais konfigūracija: vienodi svoriai $w_j = 0{,}2$ visoms penkioms atstumo funkcijoms, slenkstis $S = 0{,}5$, sezoniškumo intervalas $\pm 20$~dienų aplink tikslinę datą per visus turimus metus.
    Įvesties langas 5 valandos, o prognozės langas -- 12 valandų.

    Duomenys padalinti į mokymo (2012--2024 m.) ir testavimo (2025 m.) rinkinius.
    Iš jų sudaromi 17 valandų blokai (5 valandų įvestis ir 12 valandų \ce{KD_{2,5}} prognozė).
    Blokai sudaromi slenkančio lango principu -- einama per duomenis vienos valandos žingsniu ir kiekvienoje pozicijoje formuojamas 17 valandų iš eilės duomenų blokas.
    Iš taip suformuotų blokų atmetami tie, kuriuose yra bent vienos valandos tarpas.
    Po šio filtravimo gauti:
    \begin{itemize}
        \item Mokymo rinkinys: 61\,335 blokai.
        \item Testavimo rinkinys: 5\,630 blokai.
    \end{itemize}

    Testavimo metu kiekvieno bloko pirmosios 5 valandos sudaro įvesties langą, o paskutinės 12 -- tikrąsias reikšmes, kurios lyginamos su modelio prognoze.

    Modelio prognozavimo tikslumas vertinamas normalizuota vidutine absoliučiąja paklaida (NMAE), apibrėžta ankstesnėje metodologijos ataskaitoje:
    \begin{equation}
        NAE_t = \frac{|y_t - \hat{y}_t|}{y_{\max} - y_{\min}} \times 100\%, \qquad
        NMAE_H = \frac{1}{H} \sum_{t=1}^{H} NAE_t
    \end{equation}
    čia $y_t$ -- tikroji \ce{KD_{2,5}} reikšmė $t$-ąją prognozuojamą valandą, $\hat{y}_t$ -- modelio prognozė, $y_{\max}, y_{\min}$ -- \ce{KD_{2,5}} maksimumas ir minimumas testavimo rinkinyje, $H$ -- prognozuojamų valandų skaičius.

    Min-max normalizavimo parametrai ($x_{\min}$ ir $x_{\max}$ kiekvienam atributui) apskaičiuoti tik iš mokymo duomenų (2012--2024 m.) ir įsiminti.
    Testinio rinkinio reikšmės normalizuotos pagal tas pačias mokymo $x_{\min}$, $x_{\max}$ reikšmes, neperskaičiuojant jų iš 2025~m. duomenų.
    Tai užtikrina, kad testiniai duomenys neturi įtakos modelio prognozei, nėra duomenų nutekėjimo.

    Dėl šios priežasties testavimo rinkinio normalizuotas \ce{KD_{2,5}} rėžis nebūtinai užima visą $[-1, 1]$ intervalą.
    Faktinis testinių duomenų rėžis yra $y_{\max} - y_{\min} = 1{,}27$, kur $y_{\max} = 0{,}27$ ir $y_{\min} = -1{,}00$.
    Maksimumas $0{,}27$ gautas todėl, kad 2025 m. didžiausia \ce{KD_{2,5}} koncentracija buvo 81~µg/m\textsuperscript{3}, o istorinis maksimumas -- 127{,}60~µg/m\textsuperscript{3}, fiksuotas 2017-01-23.
    Vertinimo metrikoms apskaičiuoti naudojamas šis faktinis rėžis, o visi rezultatai pateikiami procentais nuo jo.


    \section{Bazinio KNN modelio paklaidos}

    Bazinio modelio metrikos testavimo rinkinyje pateiktos \ref{tab:baseline}~lentelėje.
    Čia $\mu_H$ -- vidutinė normalizuota absoliučioji paklaida (NMAE) $H$ valandoms, apskaičiuota per visus testinius langus, kuriems modelis pateikė prognozę, $\sigma_H$ -- tos paklaidos standartinis nuokrypis.
    Padengimas $C$ nusako kuriai daliai testinių duomenų modelis pateikė prognozę.
    \begin{table}[H]
        \centering
        \caption{Bazinio modelio metrikos ($w_j = 0{,}2$, $S = 0{,}5$)}
        \small
        \setlength{\tabcolsep}{10pt}
        \begin{tabular}{|l|r|}
            \hline
            \textbf{Metrika} & \textbf{Reikšmė} \\
            \hline
            $\mu_5$          & 5{,}87\,\%       \\
            $\sigma_5$       & 4{,}24\,\%       \\
            $\mu_{12}$       & 6{,}57\,\%       \\
            $\sigma_{12}$    & 4{,}27\,\%       \\
            Padengimas $C$   & 84{,}71\,\%       \\
            \hline
        \end{tabular}
        \label{tab:baseline}
    \end{table}

    Modelis pateikė prognozę 84,71\% testinių langų. Likusiems 15,3\% atvejų istoriniame rinkinyje nerasta pakankamai panašių 5 valandų langų pagal slenksčio sąlygą $D \le S$.
    12 valandų vidutinė normalizuota absoliučioji paklaida sudaro 6,57\%, o 5 valandų -- 5,87\%.

    0,7 procentinio punkto skirtumas tarp $\mu_5$ ir $\mu_{12}$ atsiranda todėl, kad kaimynai parenkami pagal pirmas 5 valandas, o prognozuojama 12~valandų į priekį -- per ilgesnį laikotarpį meteorologinės sąlygos pasikeičia, ir kaimyno taršos eiga nutolsta nuo tikslo realios eigos.


    \section{Paklaida pagal prognozės valandą}

    Vidutinė normalizuota absoliučioji paklaida $\mu_t$ ir jos standartinis nuokrypis $\sigma_t$ apskaičiuoti atskirai kiekvienai iš 12~prognozės valandų.
    Rezultatai pateikiami \ref{tab:perhour}~lentelėje.

    \begin{table}[H]
        \centering
        \caption{Vidutinė paklaida ir standartinis nuokrypis pagal prognozės valandą}
        \small
        \setlength{\tabcolsep}{10pt}
        \begin{tabular}{|c|r|r|}
            \hline
            \textbf{Valanda $t$} & \textbf{$\mu_t$ (\%)} & \textbf{$\sigma_t$ (\%)} \\
            \hline
            1                    & 5{,}28                & 4{,}74                   \\
            2                    & 5{,}6                 & 5{,}17                   \\
            3                    & 5{,}93                & 5{,}52                   \\
            4                    & 6{,}17                & 5{,}67                   \\
            5                    & 6{,}36                & 5{,}75                   \\
            6                    & 6{,}57                & 5{,}9                    \\
            7                    & 6{,}77                & 6{,}08                   \\
            8                    & 6{,}98                & 6{,}22                   \\
            9                    & 7{,}15                & 6{,}28                   \\
            10                   & 7{,}21                & 6{,}4                    \\
            11                   & 7{,}33                & 6{,}62                   \\
            12                   & 7{,}43                & 6{,}81                   \\
            \hline
        \end{tabular}
        \label{tab:perhour}
    \end{table}

    Pirmosios valandos paklaida yra 5,28\%, dvyliktosios -- 7,43\%.
    Paklaida didėja kiekvieną valandą ir per 12 valandų išauga 2,14 procentino punkto.
    Standartinis nuokrypis taip pat auga: nuo 4,74\% pirmajai valandai iki 6,81\% dvyliktajai.
    Tai reiškia, kad tolimesnėms valandoms ne tik vidutinė paklaida didesnė, bet ir atskirų langų paklaidos labiau skiriasi viena nuo kitos.


    \section{Atvejai kai modelis nepateikia prognozės}

    Modelyje numatytas atsisakymo mechanizmas: jei istoriniame rinkinyje nerandama pakankamai panašių situacijų pagal slenksčio sąlygą $D \le S$, modelis nepateikia prognozės.
    Bazinei konfigūracijai iš 5\,630 testinių langų prognozė nepateikta 861~atveju (15,29\%).

    \begin{table}[H]
        \centering
        \caption{Atvejų, kai nėra prognozės, pasiskirstymas pagal dienos kategoriją}
        \small
        \setlength{\tabcolsep}{10pt}
        \begin{tabular}{|l|r|r|r|}
            \hline
            \textbf{Kategorija} & \textbf{Iš viso langų} & \textbf{Nėra prognozės} & \textbf{Dalis (\%)} \\
            \hline
            Darbo diena         & 2\,974                 & 359                     & 12{,}07             \\
            Poilsio diena       & 1\,693                 & 247                     & 14{,}59             \\
            Prieššventinė       & 963                    & 255                     & 26{,}48             \\
            \hline
        \end{tabular}
        \label{tab:no_pred_cat}
    \end{table}

    Prieššventinės dienos turi žymiai didesnę dalį, kai prognozės nėra (26,48\%) nei darbo (12,07\%) ar poilsio (14,59\%) dienos.
    Prieššventinės dienos sudaro apie 15\% bendro duomenų rinkinio, kandidatų rinkinys po visų filtrų (valanda, dienos kategorija, $\pm 20$~dienų sezonas) tampa maždaug 3--4~kartus mažesnis nei darbo dienų rinkinys.

    Šis atsisakymo mechanizmas yra svarbus modelio bruožas, kurio neuroniniai tinklai ir kiti juodosios dėžės metodai neturi -- jie prognozuoja visais atvejais, net kai istoriniame rinkinyje nėra pagrindo prognozei, todėl jų prognozės gali būti klaidingai užtikrintos retoms ar nepažintoms situacijoms.


    \section{Prognozių pavyzdžiai}

    Norint vizualiai pamatyti modelio elgseną įvairiais atvejais, atrenkami trys testinių langų rinkiniai pagal $\text{NMAE}_{12}$: šeši sėkmingiausi, šeši vidutiniai (apie medianą) ir šeši blogiausi atvejai.
    Kiekvieno grafiko antraštėje pateikiama testinio lango pradžios data, valanda, $\text{NMAE}_{12}$ ir slenkstį tenkinusių kaimynų skaičius.

    \begin{figure}[H]
        \centering
        \includegraphics[width=\textwidth]{images/examples_best.png}
        \caption{Šeši sėkmingiausi prognozės atvejai (mažiausias $\text{NMAE}_{12}$). Prognozės kreivė beveik sutampa su tikrąja, paklaidos siekia 1--2\%.}
        \label{fig:best_examples}
    \end{figure}

    \begin{figure}[H]
        \centering
        \includegraphics[width=\textwidth]{images/examples_typical.png}
        \caption{Šeši vidutiniai prognozės atvejai (atrinkti apie $\text{NMAE}_{12}$ medianą $\approx 5{,}4\,\%$).}
        \label{fig:typical_examples}
    \end{figure}

    \begin{figure}[H]
        \centering
        \includegraphics[width=\textwidth]{images/examples_worst.png}
        \caption{Šeši blogiausi prognozės atvejai (didžiausias $\text{NMAE}_{12}$, iki $\approx 47\,\%$).}
        \label{fig:worst_examples}
    \end{figure}

    Sėkmingiausi ir vidutiniai atvejai rodo tipinę modelio elgseną: bendras taršos lygis pagaunamas teisingai, bet staigūs šuoliai, pavyzdžiui, 2025-12-09 h11 pikas lieka neaptikti, nes prognozė yra kelių kaimynų vidurkis.
    Blogiausi atvejai yra kitokios prigimties -- visi šeši 2025~m. kovo mėnesį, modelis remiasi 1--3~kaimynais, kurių taršos lygis stipriai skiriasi nuo tikrojo.


    \section{Dienos kategorijos filtro tobulinimas}

    Modelio filtravimas pagal dienos kategoriją taikomas tik įvesties lango pradžios valandai.
    Tai garantuoja, kad istoriniai kaimynai prasideda toje pačioje dienos kategorijoje, bet negarantuoja, kad jų 12 valandų prognozė patenka į tą pačią kategoriją kaip ir tikslinio lango.

    Pavyzdžiui, jei tikslinis įvesties langas yra darbo dienos vakaras (23:00), o po jo seka prieššventinė diena, tikslinio lango prognozė iš dalies patenka į prieššventinę dieną.
    Filtravimui problemiški yra du perėjimai: darbo $\to$ prieššventinė (10,1\% testinių langų) ir poilsio $\to$ darbo (9,1\%) -- iš viso 19,2\%.
    Trečiasis kategorijos pokytis, prieššventinė $\to$ poilsio (11,5\%), yra deterministinis, nes prieššventinė pagal apibrėžimą reiškia, kad kita diena yra poilsio, todėl filtravimo trūkumo nesukuria.
    Galimas sprendimas -- į filtravimą įtraukti ne tik įvesties lango pradžios, bet ir prognozės pabaigos dienos kategoriją.


    \section{Rezultatai}

    Eksperimento metu realizuotas svertinis KNN metodas su bazine konfigūracija ($w_j = 0{,}2$, $S = 0{,}5$).
    Pagrindinės išvados:

    \begin{itemize}
        \item Bazinis modelis pasiekia vidutinę 12~valandų paklaidą 6,57\% su 84,71\% padengimu.
        Kas valandinė analizė parodė, kad paklaida didėja nuo 5,28\% pirmajai prognozės valandai iki 7,43\% dvyliktajai.

        \item Iš 5\,630 testinių langų 15,29\% atveju modelis atsisakė pateikti prognozę.
        Didžiausią dalį atvejų, kai modelis nepateikia prognozės, sudaro prieššventinės dienos (26,48\%). Tai paaiškinama mažesniu prieššventinių dienų duomenų kiekiu.
        Tai parodo baltosios dėžės principą -- modelis informuoja apie istorinio konteksto trūkumą, o ne pateikia nepatikimą prognozę.

        \item Identifikuotas dienos kategorijos filtravimo trūkumas: 19,2\% testinių langų patenka į nedeterministinius dienos kategorijos perėjimus (darbo $\to$ prieššventinė ir poilsio $\to$ darbo). Galimas sprendimas -- į filtravimą įtraukti ir prognozės pabaigos dienos kategoriją. Padengimas šiek tiek sumažės.
    \end{itemize}

\end{document}
